\section{REVIEW PROCESS AND EVIDENCE BASE}
\label{sec:review_methods}
\label{sec:corpus_appraisal_results}

With these comparison conditions established, the analysis turns to the
evidence base on which the technical synthesis rests. The corpus was assembled
through a systematic review reported in accordance with the Preferred Reporting
Items for Systematic Reviews and Meta-Analyses (PRISMA) 2020 statement, with
search reporting guided by PRISMA-S \cite{PRISMA2020,PRISMAS2021}. The following
subsections summarize corpus construction and the rules that preserve the
context of reported results throughout the synthesis.
Table~\ref{tab:evidence_reconciliation} provides an overview of how the
extracted evidence is used in the survey.

\begin{table*}[t]
\caption{Use of extracted evidence in the survey. Row counts are coding
records; study counts overlap across categories.}
\label{tab:evidence_reconciliation}
\centering
\footnotesize
\setlength{\tabcolsep}{5.0pt}
\renewcommand{\arraystretch}{1.04}
\begin{tabularx}{\textwidth}{@{}>{\raggedright\arraybackslash}p{0.22\textwidth}
>{\raggedleft\arraybackslash}p{0.08\textwidth}
>{\raggedleft\arraybackslash}p{0.09\textwidth}
>{\raggedright\arraybackslash}X@{}}
\toprule
Evidence use & Rows & Studies & Role in the survey \\
\midrule
Quantitative evidence & 4,997 & 206 & Numerical or directional synthesis under recorded conditions. \\
Qualitative evidence & 3,206 & 206 & Mechanisms, trends, and limitations. \\
Context-only claims & 31 & 21 & Interpretation only; outside the primary synthesis. \\
Conflict-restricted claims & 72 & 31 & Withheld from primary synthesis; unaffected evidence retained. \\
Primary synthesis & 8,203 & 206 & Quantitative and qualitative evidence combined. \\
\bottomrule
\end{tabularx}
\vspace{1pt}
\parbox{\textwidth}{\footnotesize\emph{Note:} The governed ledger contains
8,306 rows. The 8,203 primary synthesis rows and their source locations are
reported in Electronic Supplement S-Evidence; governance rules are reported in
S-Protocol.}
\end{table*}

\subsection{Corpus Construction}

To assemble that evidence base, we searched Scopus, IEEE Xplore, ScienceDirect,
SpringerLink, Wiley Online Library, and Taylor \& Francis Online. The search
covered peer-reviewed journal articles and full conference papers first made
publicly available from 1 January 2020 through 22 June 2026. We retained reports
that contained sufficient full technical content in English, examined
communication and sensing within an optical or photonic system, and provided
assessable evidence for at least one survey domain. Reviews and other contextual
sources helped position the field but did not enter the technical denominator.
Search strategies, eligibility decisions, and protocol documentation are
provided in Electronic Supplements S-Search, S-Exclusions, and S-Protocol.

After deduplication, screening, and full-text assessment, 227 reports met the
eligibility criteria. Twenty-one were companion reports describing studies also
represented by another eligible report. Consolidating the report set at study
level produced 206 unique studies for the technical synthesis; the difference
therefore reflects mapping rather than further exclusion. This mapping
preserves report-level provenance while counting each underlying study once.
Electronic Supplement ST-01 cites all 206 studies and traces their 227 reports.

% ============================================================================
% FUTURE FIGURE 3 PRODUCTION SPECIFICATION --- DO NOT MATERIALIZE IN THIS PASS
% Intended placement: here, after the corpus-construction summary.
%
% Working title:
% Selection of the O-ISAC Evidence Base
% Stable label: fig:prisma_report_study_flow
% Blueprint ID: FIG-OISAC-02
% Data authority: locked Phase C selection flow and final report-to-study
% overlay. Native units are records, reports, and unique studies.
%
% Required main path:
% 1,733 records identified -> 1,259 records screened -> 330 reports sought ->
% 272 full-text reports assessed -> 227 eligible reports -> 206 unique studies.
%
% Required side information:
% - 472 duplicates and 2 other records removed before screening.
% - 927 records did not advance: 864 excluded, 61 retained as contextual, and
%   2 marked as duplicate or related versions.
% - 332 reports were forwarded and 2 aliases consolidated, leaving 330 unique
%   reports sought; 58 were not retrieved and 272 were assessed.
% - Full-text assessment produced 39 exclusions, 6 contextual reports, and 227
%   eligible reports.
% - The contextual corpus contains 67 items: 61 records plus 6 reports.
% - The 21-report difference is companion-report consolidation, not exclusion.
%
% Reader question:
% How did the search results become the study corpus used by the survey?
%
% Main visual message:
% The PRISMA flow documents selection, while report-to-study consolidation
% establishes the denominator used in the technical synthesis.
%
% Future lead-in (activate only after production and audit):
% Figure~\ref{fig:prisma_report_study_flow} summarizes study selection and the
% final mapping from eligible reports to unique studies.
%
% Future caption:
% Selection of the O-ISAC evidence base. The flow moves from bibliographic
% records to eligible reports and then to unique studies after companion reports
% are consolidated.
%
% Visual constraints:
% - Keep records, reports, and studies visibly distinct.
% - Do not display the predecessor OSF snapshot, search-history aliases, claim
%   accounting, TQAF scores, or evidence-body counts in this figure.
% ============================================================================

\subsection{Evidence Extraction and Use}

With the study denominator established, we extracted evidence at three linked
levels: study, report, and individual claim. For each result, we recorded its
definition, unit, measurement plane, operating conditions, validation setting,
and exact source location. Results reported under different conditions remained
separate; missing information was not inferred, and graphs were not digitized
to create new values. Electronic Supplement S-Data Dictionary provides the
complete coding schema.

Claims extracted from an included study could serve different roles.
Quantitative and qualitative claims could support the synthesis within their
recorded context; context-only claims informed interpretation, whereas
conflict-restricted claims were withheld from the primary synthesis. Relations
across studies also had to satisfy the contextual alignment rules in
Table~\ref{tab:comparison_record}.

Because these decisions operated at claim level, restricting one claim did not
discard unaffected evidence from the same study. All 206 studies therefore
remained represented in the primary synthesis by claims supported within their
recorded context.

\subsection{Technical Appraisal and Synthesis}

To clarify how each study could support the synthesis, we characterized the
technical evidence and reporting of all 206 studies with the review-specific
Technical Quality Assessment Framework (TQAF). The framework applies
deterministic rules across eight dimensions and records overall evidence
contribution as a separate category. The resulting profile records how the
reported evidence can be interpreted in its original setting and compared
through a shared benchmark or across platforms. Overall contribution refers
specifically to the support a study provides for this synthesis. Electronic
Supplements S-Protocol and S-Appraisal provide the framework's methodological
scope, scoring rules, and study-level records.

Overall contribution was strong for 125 studies, adequate for 75, and low for
6. Across the eight dimensions, evidence was generally easier to interpret
within its reported setting than to transfer to a shared benchmark or direct
cross-platform comparison.

% ============================================================================
% FUTURE FIGURE 4 PRODUCTION SPECIFICATION --- DO NOT MATERIALIZE IN THIS PASS
% Intended placement: here, after the concise TQAF interpretation.
%
% Working title:
% Technical Appraisal Profile of the 206-Study Evidence Base
% Stable label: fig:tqaf_profile
% Blueprint ID: FIG-OISAC-03
% Data authority: final Phase E TQAF profile for all 206 unique studies.
%
% Reader question:
% Which qualities make the evidence interpretable, and which limit portable
% comparison and reuse?
%
% Layout: one horizontal 100-percent stacked bar for each of the eight TQAF
% dimensions, grouped under two neutral braces: interpretation and reporting
% (technical relevance, metric clarity, reporting completeness, limitation
% transparency) and portability and reuse (validation maturity,
% reproducibility, benchmark readiness, comparison admissibility). Add a
% separate overall evidence-contribution bar below the eight dimensions. Use
% strong, adequate, and low segments with direct counts and percentages.
%
% Required counts:
% technical relevance 123/68/15; metric clarity 168/7/31; reporting
% completeness 196/10/0; limitation transparency 153/44/9; validation maturity
% 6/168/32; reproducibility 3/199/4; benchmark readiness 0/158/48; comparison
% admissibility 4/10/192; overall contribution 125/75/6.
%
% Main visual message:
% Interpretability and reporting are generally stronger than benchmark
% portability, reproducibility, and direct comparison admissibility.
%
% Future lead-in (activate only after production and audit):
% Figure~\ref{fig:tqaf_profile} presents the eight appraisal dimensions and the
% separate overall evidence-contribution category for all 206 studies.
%
% Future caption:
% Review-specific technical appraisal profile of the 206 included studies.
% Bars show strong, adequate, and low classifications for eight dimensions and
% a separate overall contribution category.
%
% Visual constraints:
% - Use the 206-study denominator for every bar and show zero explicitly.
% - Keep overall contribution separate from the eight input dimensions.
% - Use a colorblind-safe palette with direct labels and grayscale distinction.
% - Do not mix validation tiers or evidence-body certainty into this figure.
% ============================================================================

We therefore used a structured narrative synthesis that retained each result's
native task and unit together with its operating conditions and validation
setting. This structure preserved the meaning of heterogeneous results while
allowing recurring mechanisms and design relationships to be compared across
the corpus. The evidence was organized into 115 bodies across seven survey
domains, keeping each conclusion tied to its supporting studies. Electronic
Supplement S-Bodies provides their criteria, ratings, and memberships. With
this evidence structure in place, the technical synthesis begins with the
physical role of optics and the mechanism that couples communication and
sensing.
